\reftitle{References}

\begin{thebibliography}{99}

\bibitem{ref1}
American Psychiatric Association. \textit{Diagnostic and Statistical Manual of Mental Disorders (DSM-5)} ; American Psychiatric Publishing , 2013 ; Volume 5 , pp. 1--947.

\bibitem{ref2}
Jinan Zeidan; Eric Fombonne; Julie Scorah; Alaa Ibrahim; Maureen S. Durkin; Shekhar Saxena; Afiqah Yusuf; Andy Shih; Mayada Elsabbagh. Global prevalence of autism: A systematic review update. \textit{Autism Research} \textbf{2022} \textit{15}, \textit{5} 778--790.

\bibitem{ref3}
Catherine Lord; Mayada Elsabbagh; Gillian Baird; Jeremy Veenstra-Vanderweele. Autism spectrum disorder. \textit{The Lancet} \textbf{2018} \textit{392}, \textit{10146} 508--520.

\bibitem{ref4}
Seiji Ogawa; Tso-Ming Lee; Alan R Kay; David W Tank. Brain magnetic resonance imaging with contrast dependent on blood oxygenation. \textit{Proceedings of the National Academy of Sciences} \textbf{1990} \textit{87}, \textit{24} 9868--9872.

\bibitem{ref5}
Adriana Di Martino; Chao-Gan Yan; Qingyang Li; Erin Denio; F Xavier Castellanos; Kaat Alaerts; Jeffrey S Anderson; Michal Assaf; Susan Y Bookheimer; Mirella Dapretto; et al.. The autism brain imaging data exchange: towards a large-scale evaluation of the intrinsic brain architecture in autism. \textit{Molecular Psychiatry} \textbf{2014} \textit{19} 659--667.

\bibitem{ref6}
Lucina Q Uddin; Kaustubh Supekar; Vinod Menon. Reconceptualizing functional brain connectivity in autism from a developmental perspective. \textit{Frontiers in Human Neuroscience} \textbf{2013} \textit{7} 458.

\bibitem{ref7}
Mark Plitt; Kelly A Barnes; Alex Martin. Functional connectivity classification of autism identifies highly predictive brain features but falls short of biomarker standards. \textit{NeuroImage: Clinical} \textbf{2014} \textit{7} 359--366.

\bibitem{ref8}
J F Agastinose Ronicko; J Thomas; P Thangavel; V Koneru; G Langs; J Dauwels. Diagnostic classification of autism using resting-state fMRI data improves with full correlation functional brain connectivity compared to partial correlation. \textit{Journal of Neuroscience Methods} \textbf{2020} \textit{345} 108884.

\bibitem{ref9}
Anibal S Heinsfeld; Alexandre R Franco; R Cameron Craddock; Augusto Buchweitz; Felipe Meneguzzi. Identification of autism spectrum disorder using deep learning and the ABIDE dataset. \textit{NeuroImage: Clinical} \textbf{2017} \textit{17} 16--23.

\bibitem{ref10}
R Matthew Hutchison; Thilo Womelsdorf; Elena A Allen; Peter A Bandettini; Vince D Calhoun; Maurizio Corbetta; Stefania Della Penna; Jeff H Duyn; Gary H Glover; Javier Gonzalez-Castillo; et al.. Dynamic functional connectivity: promise, issues, and interpretations. \textit{NeuroImage} \textbf{2013} \textit{80} 360--378.

\bibitem{ref11}
Vince D Calhoun; Robyn Miller; Godfrey Pearlson; Tülay Adalı. The chronnectome: time-varying connectivity networks as the next frontier in fMRI data discovery. \textit{Neuron} \textbf{2014} \textit{84}, \textit{2} 262--274.

\bibitem{ref12}
Maria Giulia Preti; Thomas A Bolton; Dimitri Van De Ville. The dynamic functional connectome: State-of-the-art and perspectives. \textit{NeuroImage} \textbf{2017} \textit{160} 41--54.

\bibitem{ref13}
Ashish Vaswani; Noam Shazeer; Niki Parmar; Jakob Uszkoreit; Llion Jones; Aidan N Gomez; L ukasz Kaiser; Illia Polosukhin. Attention is All you Need. In \textit{Advances in Neural Information Processing Systems} , 2017 ; Volume 30.

\bibitem{ref14}
George Zerveas; Srideepika Jayaraman; Dhaval Patel; Anuradha Bhamidipaty; Carsten Eickhoff. A Transformer-based Framework for Multivariate Time Series Representation Learning. In Proceedings of the 27th ACM SIGKDD Conference on Knowledge Discovery and Data Mining , 2021 ; pp. 2114--2124.

\bibitem{ref15}
Yinchi Zhou; Peiyu Duan; Yuexi Du; Nicha C Dvornek. Self-supervised Pre-training Tasks for an fMRI Time-Series Transformer in Autism Detection. In \textit{Machine Learning in Clinical Neuroimaging} , Cham , 2025 ; pp. 145--154.

\bibitem{ref16}
Lucas Mahler; Qi Wang; Julius Steiglechner; Florian Birk; Samuel Heczko; Klaus Scheffler; Gabriele Lohmann. Pretraining is All You Need: A Multi-Atlas Enhanced Transformer Framework for Autism Spectrum Disorder Classification. In \textit{Machine Learning in Clinical Neuroimaging} , Cham , 2023 ; pp. 123--132.

\bibitem{ref17}
Xian-An Bi; Yong Wang; Qiang Shu; Qiang Sun; Qiang Xu. Classification of Autism Spectrum Disorder Using Random Support Vector Machine Cluster. \textit{Frontiers in Genetics} \textbf{2018} \textit{9} 18.

\bibitem{ref18}
Chengyu Guo; Fei Chen; Yajie Chang; Jinting Yan. Applying Random Forest classification to diagnose autism using acoustical voice-quality parameters during lexical tone production. \textit{Biomedical Signal Processing and Control} \textbf{2022} \textit{77} 103811.

\bibitem{ref19}
Huda Akil; Maryann E Martone; David C Van Essen. Challenges and opportunities in mining neuroscience data. \textit{Science} \textbf{2011} \textit{331}, \textit{6018} 708--712.

\bibitem{ref20}
Sergey M Plis; Devon R Hjelm; Ruslan Salakhutdinov; Elena A Allen; Henry J Bockholt; Jeffrey D Long; Hans J Johnson; Jane S Paulsen; Jessica A Turner; Vince D Calhoun. Deep learning for neuroimaging: a validation study. \textit{Frontiers in Neuroscience} \textbf{2014} \textit{8} 229.

\bibitem{ref21}
Alaa Bessadok; Mohamed Ali Mahjoub; Islem Rekik. Graph Neural Networks in Network Neuroscience. \textit{IEEE Transactions on Pattern Analysis and Machine Intelligence} \textbf{2023} \textit{45}, \textit{5} 5833--5848.

\bibitem{ref22}
Xiaoxiao Li; Yuan Zhou; Nicha Dvornek; Muhan Zhang; Juntang Zhuang; Pam Ventola; James Duncan. Pooling Regularized Graph Neural Network for fMRI Biomarker Analysis. In \textit{Medical Image Computing and Computer Assisted Intervention--MICCAI 2020} , 2020 ; pp. 625--635.

\bibitem{ref23}
Bryan Lim; Sercan Ö Arik; Nicolas Loeff; Tomas Pfister. Temporal Fusion Transformers for interpretable multi-horizon time series forecasting. \textit{International Journal of Forecasting} \textbf{2021} \textit{37}, \textit{4} 1748--1764.

\bibitem{ref24}
Veenu Rani; Munish Kumar; Aastha Gupta; Monika Sachdeva; Ajay Mittal; Krishan Kumar. Self-supervised learning for medical image analysis: a comprehensive review. \textit{Evolving Systems} \textbf{2024} \textit{15}, \textit{4} 1607--1633.

\bibitem{ref25}
Xiao Liu; Fanjin Zhang; Zhenyu Hou; Li Mian; Zhaoyu Wang; Jing Zhang; Jie Tang. Self-Supervised Learning: Generative or Contrastive. \textit{IEEE Transactions on Knowledge and Data Engineering} \textbf{2023} \textit{35}, \textit{1} 857--876.

\bibitem{ref26}
Yen-Wei Chen; Lanfen Lin; Rahul Kumar Jain. Multimodal Deep Learning for Medical Image Analysis. In \textit{Recent Advances in Deep Learning for Medical Image Analysis: Paradigms and Applications} ; Springer Nature Switzerland: Cham , 2026 ; pp. 129--150.

\bibitem{ref27}
R O'Halloran; Brian H Kopell; Emma Sprooten; Wayne K Goodman; Sophia Frangou. Multimodal Neuroimaging-Informed Clinical Applications in Neuropsychiatric Disorders. \textit{Frontiers in Psychiatry} \textbf{2016} \textit{7} 63.

\bibitem{ref28}
Y Cheng; L Liu; X Gu; Z Lu; Y Xia; J Chen; L Tang. Graph fusion prediction of autism based on attentional mechanisms. \textit{Journal of Biomedical Informatics} \textbf{2023} \textit{146} 104484.

\bibitem{ref29}
Swathi Pothuganti. Analysis on Solutions for Over-fitting and Under-fitting in Machine Learning Algorithms. \textit{International Journal of Innovative Research in Science Engineering and Technology} \textbf{2018} \textit{7} 12401--12404.

\bibitem{ref30}
Arthur Le Guennec; Simon Malinowski; Romain Tavenard. Data Augmentation for Time Series Classification using Convolutional Neural Networks. In ECML/PKDD Workshop on Advanced Analytics and Learning on Temporal Data , Riva Del Garda, Italy , 2016 .

\bibitem{ref31}
Jason Yosinski; Jeff Clune; Yoshua Bengio; Hod Lipson. How transferable are features in deep neural networks?. In \textit{Advances in Neural Information Processing Systems} , 2014 ; Volume 27.

\bibitem{ref32}
A. Di Martino; C.-G. Yan; Q. Li; E. Denio; F. X. Castellanos; K. Alaerts; J. S. Anderson; M. Assaf; S. Y. Bookheimer; M. Dapretto; et al.. The autism brain imaging data exchange: towards a large-scale evaluation of the intrinsic brain architecture in autism. \textit{Molecular psychiatry} \textbf{2014} \textit{19}, \textit{6} 659--667.

\bibitem{ref33}
C. Craddock; Y. Benhajali; C. Chu; F. Chouinard; A. Evans; A. Jakab; B. S. Khundrakpam; J. D. Lewis; Q. Li; M. Milham; et al.. The neuro bureau preprocessing initiative: open sharing of preprocessed neuroimaging data and derivatives. \textit{Frontiers in Neuroinformatics} \textbf{2013} \textit{7}, \textit{27} 5.

\end{thebibliography}